\documentclass[aspectratio=169]{beamer}
\usetheme{Madrid}
\usecolortheme{beaver}
\usepackage[utf8]{inputenc}
\usepackage[T1]{fontenc}
\usepackage{tikz}
\usepackage{amssymb}
\usepackage{graphicx}
\setbeamertemplate{navigation symbols}{}
\setbeamercolor{title}{fg=white,bg=blue!55!black}
\setbeamercolor{frametitle}{fg=white,bg=blue!45!black}
\setbeamercolor{block title}{bg=blue!15,fg=black}
\setbeamercolor{block body}{bg=blue!5,fg=black}

\title{Web Story 5 — Your Story Engine}
\subtitle{Everything from Semester 1, working together in one real interactive page}
\author{ByteForge --- Grade 8 --- App Builder}
\date{}

\begin{document}

\begin{frame}
    \titlepage
\end{frame}

\begin{frame}{Today's Plan}
    \begin{itemize}
        \item Add and Re-Render
        \item Track a Flag, Too
        \item Save As You Go
        \item[$\bigstar$] BOSS --- Your Story Engine
    \end{itemize}
\end{frame}

\begin{frame}{A click that triggers a re-render}
    \Large
    \texttt{addEventListener("click", ...)} is genuine JavaScript event handling. Clearing the list before looping again means the display always reflects the CURRENT inventory.
\end{frame}


\begin{frame}{One click, multiple kinds of state updated}
    \Large
    A single click genuinely updates two DIFFERENT shapes of state at once — an ordered list (inventory) and a named fact (\texttt{flags.hasVisited}).
\end{frame}


\begin{frame}{Saving on every meaningful change}
    \Large
    Real progress-saving usually happens repeatedly, every time something changes — every click that changes inventory also saves it, so nothing is ever lost.
\end{frame}


\begin{frame}{BOSS — a real, small, working story engine}
    \Large
    State, functions, loops, real click events, and persistence — this is genuinely the same shape as a real, if small, interactive web application.
\end{frame}

\begin{frame}{Real-World Example}
    \Large
    This is genuinely the shape of a real single-page app: state changes, the UI re-renders to match, and changes are saved as they happen — the exact architecture behind most modern web and mobile apps, just at a much larger scale.
\end{frame}


\begin{frame}{Let Them Discover}
    \Large
    Before opening the platform: sketch what SHOULD happen, in order, when a player clicks "pick up sword" — which piece of state changes first, what re-renders, and when does it save?
\end{frame}

\begin{frame}{Your Turn}
    \Large
    Open the student map, find \textbf{Web Story 5 — Your Story Engine}, and work through it
    hands-on --- everything you just saw, they now do themselves.
\end{frame}

\end{document}
